% !TeX root = ../main.tex
%%% ---------------------------------------------------------------------------
%%% 备用页
%%%
%%% main.tex 里的 \appendix 已经让这些帧不计入页脚的总页数，
%%% 所以听众看到的 "12/18" 不会把备用页算进去。
%%%
%%% 放什么：预判会被问到的问题的答案——完整推导、完整对比表、
%%% 失败案例、超参数、与某篇相关工作的区别。
%%% 提问环节直接翻过来，比现场画图有说服力得多。
%%% ---------------------------------------------------------------------------

\begin{frame}[plain, noframenumbering]
  \vfill
  \centering
  {\usebeamerfont{section page}\usebeamercolor[fg]{section page} 备用页\par}
  \vskip 0.9em
  {\usebeamercolor[fg]{structure}\rule{0.14\paperwidth}{0.8pt}}
  \vfill
\end{frame}

\begin{frame}{备用：方差上界的紧性}
  对 $g \in (0,1)$，由 $\Var[\featmap_l'] = g^2\sigma_1^2 + (1-g)^2\sigma_2^2$
  进一步可得下界
  \begin{equation*}
    \Var\!\left[\featmap_l'\right]
      \ge \frac{\sigma_1^2 \sigma_2^2}{\sigma_1^2 + \sigma_2^2},
  \end{equation*}
  等号在 $g^{*} = \sigma_2^2 / (\sigma_1^2 + \sigma_2^2)$ 处取得。

  \vspace{0.8em}
  \begin{block}{这说明什么}
    门控不仅不放大方差，在中间取值处实际起到了\alert{方差收缩}的作用——
    这解释了为什么加入门控后训练曲线更平滑。
  \end{block}
\end{frame}

\begin{frame}{备用：完整超参数}
  \begin{table}
    \centering\small
    \begin{tabular}{l l}
      \toprule
      {超参数} & {取值} \\
      \midrule
      优化器             & SGD（动量 \num{0.9}） \\
      初始学习率         & \num{0.01}，余弦退火，预热 \num{500} 步 \\
      权重衰减           & \num{1e-4} \\
      批大小             & 16（4 卡 A100，每卡 4） \\
      训练轮数           & \qty{36}{\epoch} \\
      输入尺寸           & $1024 \times 1024$，重叠 \qty{200}{\pixel} \\
      门控压缩比 $r$     & 16 \\
      损失权重 $\lambda$ & \num{1.0} \\
      \bottomrule
    \end{tabular}
  \end{table}
\end{frame}

\begin{frame}{备用：失败案例}
  \begin{columns}[T, onlytextwidth]
    \begin{column}{0.46\textwidth}
      \centering
      \placeholder{\linewidth}{3.8cm}{figures/failure-case.pdf}
    \end{column}
    \begin{column}{0.50\textwidth}
      \begin{itemize}
        \setlength{\itemsep}{0.7em}
        \item 云层遮挡下门控权重失去区分度，退化为均匀融合
        \item 极端长宽比目标（如跑道）仍存在定位偏差
        \item 这两类场景在测试集中占 \qty{4.1}{\percent}
      \end{itemize}
    \end{column}
  \end{columns}
\end{frame}
